% GENERATED FILE. Edit canonical lessons, then run npm run generate:books.
% canonical-source-hash: fnv1a64:eb7a10a1b9262c3a
% canonical-lessons: GE-C363-umweltschutz, GE-C363-klimawandel, GE-C363-nachhaltig, GE-C363-recyceln, GE-C363-verbrauchen

\chapter{Die Umwelt schützen}
\label{ch:ge-die-umwelt-schutzen}
\hlchaptermodality{363}

\begin{chapteropening}
\textbf{By the end of this chapter:} \emph{I can talk about protecting the environment, climate change, sustainability, recycling and consumption.}

Complete the last lesson of chapter 363: I can talk about protecting the environment, climate change, sustainability, recycling and consumption.
\end{chapteropening}

\section[der Umweltschutz]{der Umweltschutz — environmental protection}

\label{lesson:GE-C363-umweltschutz}



\begin{quote}
Before the new one: say the German for to donate, then the German for to get involved.
\end{quote}

\subsection*{You'll want to know: der Umweltschutz}
\textbf{der Umweltschutz} — \textquotedblleft{}environmental protection\textquotedblright{}.

\begin{quote}
Für den Umweltschutz fahren wir mit dem Fahrrad zur Arbeit.
\end{quote}

\begin{quote}
\emph{We cycle to work to protect the environment.}
\end{quote}

No plural. Built from \emph{die Umwelt} — the environment — and \emph{der Schutz} — protection.

\begin{cousinweb}
\emph{Umweltschutz} joins \emph{Umwelt} \textquotedblleft{}environment\textquotedblright{}, from \emph{um} \textquotedblleft{}around\textquotedblright{} and \emph{Welt} \textquotedblleft{}world\textquotedblright{}, Middle High German \emph{welt} (EtymDB), to \emph{Schutz} \textquotedblleft{}protection\textquotedblright{}, Middle High German \emph{schuz} (EtymDB). \emph{Umwelt} is usually said to render Danish \emph{omverden} \textquotedblleft{}the world around\textquotedblright{}.
\end{cousinweb}

\begin{grammarlens}[title={What you've built}]
One more word for telling, explaining and describing at B1.
\end{grammarlens}

\subsection*{Your turn}
\begin{itemize}
\raggedright
  \item \emph{Say it:} \emph{der Umweltschutz}
  \item \emph{Say it:} \emph{der Umweltschutz}, once more
  \item \emph{Say it:} a sentence of your own with \emph{der Umweltschutz}
  \item \emph{Recall:} say the German for to donate, then the German for to get involved, then say \emph{der Umweltschutz} again
\end{itemize}

\begin{tcolorbox}[breakable,colback=teal!4,colframe=teal!35!black,title={Before you move on}]
What does der Umweltschutz mean? (\textquotedblleft{}Environmental protection\textquotedblright{}.) Say it once more.
\end{tcolorbox}
\section[der Klimawandel]{der Klimawandel — climate change}

\label{lesson:GE-C363-klimawandel}



\begin{quote}
Before the new one: say the German for to get involved, then the German for environmental protection.
\end{quote}

\subsection*{You'll want to know: der Klimawandel}
\textbf{der Klimawandel} — \textquotedblleft{}climate change\textquotedblright{}.

\begin{quote}
Der Klimawandel macht die Sommer immer heißer.
\end{quote}

\begin{quote}
\emph{Climate change is making the summers hotter and hotter.}
\end{quote}

No plural. \emph{Das Klima} — the climate — and \emph{der Wandel} — change.

\begin{cousinweb}
\emph{Klimawandel} joins \emph{Klima}, from Latin \emph{clima} and Greek \emph{klíma} \textquotedblleft{}slope, zone of latitude\textquotedblright{}, from \emph{klínein} \textquotedblleft{}to slope\textquotedblright{} (EtymDB), to \emph{Wandel} \textquotedblleft{}change\textquotedblright{}, inherited from Old High German \emph{wandil} (EtymDB).
\end{cousinweb}

\begin{grammarlens}[title={What you've built}]
One more word for telling, explaining and describing at B1.
\end{grammarlens}

\subsection*{Your turn}
\begin{itemize}
\raggedright
  \item \emph{Say it:} \emph{der Klimawandel}
  \item \emph{Say it:} \emph{der Klimawandel}, once more
  \item \emph{Say it:} a sentence of your own with \emph{der Klimawandel}
  \item \emph{Recall:} say the German for to get involved, then the German for environmental protection, then say \emph{der Klimawandel} again
\end{itemize}

\begin{tcolorbox}[breakable,colback=teal!4,colframe=teal!35!black,title={Before you move on}]
What does der Klimawandel mean? (\textquotedblleft{}Climate change\textquotedblright{}.) Say it once more.
\end{tcolorbox}
\section[nachhaltig]{nachhaltig — sustainable; lasting}

\label{lesson:GE-C363-nachhaltig}



\begin{quote}
Before the new one: say the German for environmental protection, then the German for climate change.
\end{quote}

\subsection*{You'll want to know: nachhaltig}
\textbf{nachhaltig} — \textquotedblleft{}sustainable; lasting\textquotedblright{}.

\begin{quote}
Wir kaufen nur noch nachhaltige Produkte.
\end{quote}

\begin{quote}
\emph{We only buy sustainable products now.}
\end{quote}

\emph{Die Nachhaltigkeit} — sustainability. \emph{Ein nachhaltiger Eindruck} — a lasting impression.

\begin{cousinweb}
\emph{nachhaltig} is built on \emph{Nachhalt} \textquotedblleft{}something held back, a reserve\textquotedblright{}, from \emph{nach} and \emph{halten} \textquotedblleft{}to hold\textquotedblright{}, Proto-Germanic \emph{haldaną} (EtymDB). It is usually said to have been first a word of forestry, for felling no more trees than grow back.
\end{cousinweb}

\begin{grammarlens}[title={What you've built}]
One more word for telling, explaining and describing at B1.
\end{grammarlens}

\subsection*{Your turn}
\begin{itemize}
\raggedright
  \item \emph{Say it:} \emph{nachhaltig}
  \item \emph{Say it:} \emph{nachhaltig}, once more
  \item \emph{Say it:} a sentence of your own with \emph{nachhaltig}
  \item \emph{Recall:} say the German for environmental protection, then the German for climate change, then say \emph{nachhaltig} again
\end{itemize}

\begin{tcolorbox}[breakable,colback=teal!4,colframe=teal!35!black,title={Before you move on}]
What does nachhaltig mean? (\textquotedblleft{}Sustainable; lasting\textquotedblright{}.) Say it once more.
\end{tcolorbox}
\section[recyceln]{recyceln — to recycle}

\label{lesson:GE-C363-recyceln}



\begin{quote}
Before the new one: say the German for climate change, then the German for sustainable.
\end{quote}

\subsection*{You'll want to know: recyceln}
\textbf{recyceln} — \textquotedblleft{}to recycle\textquotedblright{}.

\begin{quote}
Wir recyceln Papier, Glas und Plastik.
\end{quote}

\begin{quote}
\emph{We recycle paper, glass and plastic.}
\end{quote}

Said the English way: \emph{ri-SAI-keln}. Perfekt: \emph{hat recycelt}.

\begin{cousinweb}
\emph{recyceln} was borrowed from English \emph{recycle} (EtymDB), built from \emph{re-} \textquotedblleft{}again\textquotedblright{} and \emph{cycle} \textquotedblleft{}a round\textquotedblright{}, which goes back through Latin \emph{cyclus} to Greek \emph{kýklos} \textquotedblleft{}a ring\textquotedblright{}.
\end{cousinweb}

\begin{grammarlens}[title={What you've built}]
One more word for telling, explaining and describing at B1.
\end{grammarlens}

\subsection*{Your turn}
\begin{itemize}
\raggedright
  \item \emph{Say it:} \emph{recyceln}
  \item \emph{Say it:} \emph{recyceln}, once more
  \item \emph{Say it:} a sentence of your own with \emph{recyceln}
  \item \emph{Recall:} say the German for climate change, then the German for sustainable, then say \emph{recyceln} again
\end{itemize}

\begin{tcolorbox}[breakable,colback=teal!4,colframe=teal!35!black,title={Before you move on}]
What does recyceln mean? (\textquotedblleft{}To recycle\textquotedblright{}.) Say it once more.
\end{tcolorbox}
\section[verbrauchen]{verbrauchen — to consume, to use up}

\label{lesson:GE-C363-verbrauchen}



\begin{quote}
Before the new one: say the German for sustainable, then the German for to recycle.
\end{quote}

\subsection*{You'll want to know: verbrauchen}
\textbf{verbrauchen} — \textquotedblleft{}to consume, to use up\textquotedblright{}.

\begin{quote}
Unser alter Kühlschrank verbraucht zu viel Strom.
\end{quote}

\begin{quote}
\emph{Our old fridge uses too much electricity.}
\end{quote}

Inseparable: \emph{er verbraucht, er hat verbraucht}. \emph{Der Verbrauch} — consumption.

\begin{cousinweb}
\emph{verbrauchen} is built with \emph{ver-} on \emph{brauchen} \textquotedblleft{}to need, to use\textquotedblright{}, inherited from Old High German \emph{brūhhan} and Proto-Germanic \emph{brūkaną} \textquotedblleft{}to use, to consume\textquotedblright{} (EtymDB).
\end{cousinweb}

\begin{grammarlens}[title={What you've built}]
One more word for telling, explaining and describing at B1.
\end{grammarlens}

\subsection*{Your turn}
\begin{itemize}
\raggedright
  \item \emph{Say it:} \emph{verbrauchen}
  \item \emph{Say it:} \emph{verbrauchen}, once more
  \item \emph{Say it:} a sentence of your own with \emph{verbrauchen}
  \item \emph{Recall:} say the German for sustainable, then the German for to recycle, then say \emph{verbrauchen} again
\end{itemize}

\begin{tcolorbox}[breakable,colback=teal!4,colframe=teal!35!black,title={Before you move on}]
What does verbrauchen mean? (\textquotedblleft{}To consume, to use up\textquotedblright{}.) Say it once more.
\end{tcolorbox}
